% =====================================================================
%  c00.tex —— 标题页 / 摘要 / 译者说明（Lead 独占）
% =====================================================================
\begin{titlepage}
\centering
\vspace*{2.2cm}

{\small IPMU-16-0048}\\[2.4cm]

{\LARGE\bfseries 2d/4d 对应关系简述}\\[0.6cm]
{\large\itshape A brief review of the 2d/4d correspondences}\\[2.2cm]

{\large 立川裕二（Yuji Tachikawa）}\\[0.35cm]
{\small 东京大学 卡弗里宇宙物理与数学研究所（Kavli IPMU,\\ University of Tokyo, Kashiwa, Chiba 277-8583, Japan）}\\[0.15cm]
{\small\texttt{[邮箱隐去]}}\\[1.8cm]

{\small 本文为综述文集《Localization techniques in quantum field theories》\\
（编者：V. Pestun、M. Zabzine，共 17 章）之一章。}\\[0.2cm]
{\small arXiv:1608.02964v5 [hep-th]，2016 年 10 月 15 日}\\[2.6cm]

{\small 中文译注本 \quad 译者批注以 \texttt{\% [Tutor Note: ...]} 形式\\
存放于 \LaTeX{} 源码注释中，不进入排版}
\vfill
\end{titlepage}

% =====================================================================
\section*{摘要}
\addcontentsline{toc}{section}{摘要}
% =====================================================================
本文给出 2d/4d 对应关系的一个初等导引。我们先简要回顾二维 $q$-形变杨--米尔斯理论与
Liouville 理论，然后引入把三基表示（trifundamental）耦合到 $\SU(2)$ 规范场而得到的四维理论。
随后我们具体看到：这些理论在 $S^3\times S^1$ 与 $S^4$ 上的超对称配分函数分别由 $q$-形变
杨--米尔斯理论与 Liouville 理论给出。在从六维 $\cN=(2,0)$ 理论的角度讨论这一对应关系可能
如何被理解之后，我们列举未来方向作为综述的结束。大部分技术细节将指向更详细的综述文章。

\medskip
\noindent\textit{（原摘要见原文 p.1；本译本保留原摘要的完整信息，译者说明见下。）}

% =====================================================================
\section*{译者说明}
\addcontentsline{toc}{section}{译者说明}
% =====================================================================
\begin{enumerate}[leftmargin=2em,itemsep=2pt]
\item \textbf{底本。}本译本以 arXiv:1608.02964v5（2016-10-15）为底本，共 46 页：正文
  pp.1--41（含致谢），参考文献 pp.41--46（77 条）。全文无附录、无表格，含 13 幅矢量图、
  137 条编号公式。
\item \textbf{页码映射。}底本 PDF 页码与印刷页码一致（偏移 0），故文中批注一律写作
  「原文 p.$N$」。本译本自身的页码见目录；公式编号 $(x.y)$、图号、文献编号一律照抄原文，
  与原文一一对应，可以据此回查原 PDF。
\item \textbf{体例。}章节编号与原文章节自动对应；公式一律用 \verb|\tag{X.Y}| 复刻原编号；
  图表编号手工给定，与原文一致；文献引用写作纯文本 \verb|[12]|。
\item \textbf{批注。}凡属译者补充的解释、易错点提醒、以及原文笔误的指认，一律写在
  \LaTeX{} 源码的 \verb|% [Tutor Note: ...]| 注释中，\textbf{不进入排版}；
  它们是源码的一部分，汇总见随交付的《译注汇总.md》。
\item \textbf{原文笔误。}正文一律照原文排印，译者不作径改；凡确认的笔误在批注中指出并给出
  正确形式，汇总见 README 的「原文笔误登记表」。
\item \textbf{本文性质。}这是一篇\textbf{导引性综述}，作者本人在正文中多次把技术细节推给
  更详细的综述（见文献表）。因此本译本对「作者省略的论证」不作补写式改写，只在批注中
  指出该处省略了什么、可去哪里补。
\end{enumerate}

\clearpage
